\begin{helper}[Fixed-Set External Product Model Ingredients]
Fix \(n\in\mathbb N\), a finite set \(I\subseteq\mathrm{Fin}(n)\), parameters
\(\varepsilon,\varepsilon_1,\varepsilon_2\in\mathbb R\), and \(n_0\in\mathbb N\).
Assume
\[
\noderef{ParameterHierarchyEventualComparisons}
  (\varepsilon,\varepsilon_1,\varepsilon_2,n_0),\qquad n_0<n,
\]
and
\[
|I|=\noderef{TwoBiteNaturalIndependenceScale}(1+\varepsilon,n).
\]
Put
\[
M=\noderef{TwoBiteNaturalReducedVertexCount}(n),\qquad
p_0=\noderef{TwoBiteEdgeProbability}\!\left(\frac12,n\right),
\]
\[
L=\log n,\qquad
K=\noderef{TwoBiteNaturalIndependenceScale}(1+\varepsilon,n),
\qquad
C=\noderef{RealChooseTwo}(K),
\]
where \(K\) is viewed as a real number in the formula for \(C\).  Set
\[
r_{\mathrm{prod}}=2M,\qquad
c_{\mathrm{prod}}=\frac12 n^{4\varepsilon},\qquad
t=\frac{C}{\sqrt L}.
\]
For \(\omega\in\noderef{TwoBiteSample}(n,M,p_0)\), form
\[
P_R=\{\noderef{TwoBiteRedProjection}(\omega,u):u\in I\},\qquad
P_B=\{\noderef{TwoBiteBlueProjection}(\omega,u):u\in I\},
\]
let \(P_I\) be the union of \(P_R\) in the red copy and \(P_B\) in the blue
copy, and set
\[
S_I^{\mathrm{ext}}
=\{x:\noderef{TwoBiteSmallClass}(\omega,\varepsilon,I,x)\}\setminus P_I.
\]
Define
\[
E_{\mathrm{ext}}(\omega)
=
\bigcup_{x\in S_I^{\mathrm{ext}}}
  \noderef{TwoBiteFinalPairs}(\noderef{TwoBiteX}(\omega,I,x)),
\]
and let \(Z_I(\omega)\) be the number of pairs \(e=(u,v)\) in
\(E_{\mathrm{ext}}(\omega)\) satisfying
\[
\noderef{TwoBiteRedProjection}(\omega,u)
  \ne\noderef{TwoBiteRedProjection}(\omega,v),\qquad
\noderef{TwoBiteBlueProjection}(\omega,u)
  \ne\noderef{TwoBiteBlueProjection}(\omega,v).
\]
Let
\[
\mu(\omega)=\frac{C}{2L}.
\]
For every fixed embedding
\[
\pi:\mathrm{Fin}(n)\hookrightarrow\mathrm{Fin}(M)\times\mathrm{Fin}(M),
\]
there are a finite type \(\alpha\), a probability mass \(q:\alpha\to\mathbb R\),
and a product variable
\[
Z_{\mathrm{prod}}:(\mathrm{Fin}(r_{\mathrm{prod}})\to\alpha)\to\mathbb R
\]
such that
\[
q(a)\ge0,\qquad \sum_{a\in\alpha}q(a)=1,
\]
\[
0<c_{\mathrm{prod}},\qquad 0\le t,
\]
\[
\left|Z_{\mathrm{prod}}(v)-Z_{\mathrm{prod}}(v')\right|
\le c_{\mathrm{prod}}
\]
whenever \(v\) and \(v'\) differ in at most one product coordinate,
\[
\mathbb E_q Z_{\mathrm{prod}}\le \frac{C}{2L},
\]
and
\[
\begin{aligned}
&\noderef{TwoBiteEventProbability}(n,M,p_0)
 \bigl(\{\omega:\noderef{TwoBiteEmbedding}(\omega)=\pi
   \text{ and } Z_I(\omega)\ge \mu(\omega)+t\}\bigr)\\
&\qquad\le
\Pr_q\!\left(Z_{\mathrm{prod}}\ge \mathbb E_q Z_{\mathrm{prod}}+t\right)
\noderef{TwoBiteEventProbability}(n,M,p_0)
 \bigl(\{\omega:\noderef{TwoBiteEmbedding}(\omega)=\pi\}\bigr).
\end{aligned}
\]
Here
\[
\mathbb E_q Z_{\mathrm{prod}}
=
\sum_v\left(\prod_iq(v_i)\right)Z_{\mathrm{prod}}(v),
\]
and the displayed product-tail probability is the corresponding finite sum of
\(\prod_i q(v_i)\) over the product outcomes satisfying the displayed tail
inequality.
\end{helper}
\begin{proof}
Fix \(\pi\).  On the embedding cell
\[
\mathcal C_\pi=\{\omega:\noderef{TwoBiteEmbedding}(\omega)=\pi\},
\]
the red and blue projections of every \(u\in I\) are fixed: they are
\(\pi(u)_1\) and \(\pi(u)_2\).  Thus the sets \(P_R\) and \(P_B\), and hence
the embedded set \(P_I\), are fixed on the cell.  If
\(\noderef{TwoBiteEventProbability}(\mathcal C_\pi)=0\), the final probability
inequality is immediate because the event on the left is a subevent of
\(\mathcal C_\pi\).  Assume from now on that the cell has positive mass.

Unfolding \(\noderef{TwoBiteEventProbability}\), the weight of a sample is
\[
\noderef{GnpGraphWeight}(p_0,G_R)\,
\noderef{GnpGraphWeight}(p_0,G_B)\,
\noderef{UniformInjectionWeight}(n,M).
\]
After conditioning on \(\mathcal C_\pi\), the injection factor is constant and
cancels.  Therefore the conditional law of the two base graphs is the product
of the red and blue \(G(M,p_0)\) graph laws.  Any edge coordinate not incident
to \(P_R\) in the red graph or to \(P_B\) in the blue graph can be summed out;
for such an unused coordinate the two possible values contribute
\((1-p_0)+p_0=1\).

Let
\[
\alpha=\mathcal P(\mathrm{Fin}(M))\times\mathcal P(\mathrm{Fin}(M)).
\]
For \(a=(A_R,A_B)\in\alpha\), define \(q(a)\) to be
\[
\noderef{FixedSetProductModelMassFn}(p_0,P_R,P_B)(a).
\]
Equivalently,
\[
q(a)=p_0^{|A_R|}(1-p_0)^{|P_R|-|A_R|}
     \cdot p_0^{|A_B|}(1-p_0)^{|P_B|-|A_B|}
\]
when \(A_R\subseteq P_R\) and \(A_B\subseteq P_B\), and \(q(a)=0\)
otherwise.  The hierarchy gives \(0\le p_0\le1\), so
\(\noderef{FixedSetProductModelMass}\) applies and yields
\[
q(a)\ge0\quad\text{for all }a,\qquad \sum_aq(a)=1.
\]

Use \(2M\) product coordinates.  The first \(M\) coordinates are indexed by
red base vertices and use only the first component \(A_R\); the last \(M\)
coordinates are indexed by blue base vertices and use only the second
component \(A_B\).  Coordinates whose chosen \(A_R\) or \(A_B\) is not
contained in \(P_R\) or \(P_B\) have \(q\)-mass zero.  If a coordinate
corresponds to a vertex already in \(P_R\) or \(P_B\), the eventual
\(Z_I\)-calculation ignores it because \(P_I\) is removed.  The unused
component of each \(\alpha\)-coordinate factors through the mass-one
calculation in \(\noderef{FixedSetProductModelMass}\), so this product model
has exactly the same marginal law as the conditional law of the external
adjacency vectors that determine \(Z_I\).

Define \(Z_{\mathrm{prod}}\) by reconstructing \(Z_I\) from these external
adjacency vectors and the fixed embedding \(\pi\).  A red coordinate with
first component \(A_R\) contributes the fiber
\[
\{u\in I:\pi(u)_1\in A_R\},
\]
and a blue coordinate with second component \(A_B\) contributes
\[
\{u\in I:\pi(u)_2\in A_B\}.
\]
The small-class test, the final-pair set, and the distinct-projection filter
are then exactly the same deterministic operations used in the definition of
\(Z_I\).  Therefore the product distribution of \(Z_{\mathrm{prod}}\) is the
conditional distribution of \(Z_I\) on \(\mathcal C_\pi\).

We next compare the product expectation with \(C/(2L)\).  Fix a rank-oriented
pair \((u,w)\) from a two-element subset of \(I\) whose red and blue
projections are both distinct.
For a red external base vertex \(x\) to make this pair contribute, the red
edges from \(x\) to \(\pi(u)_1\) and \(\pi(w)_1\) must both be present.  These
are distinct non-loop coordinates under the fixed embedding, so their joint
probability is \(p_0^2\).  Requiring \(x\) to be in the small class can only
decrease the probability.  There are at most \(M\) red choices for \(x\), so
the red contribution probability is at most \(Mp_0^2\).  The same argument in
the blue graph gives another \(Mp_0^2\).  Hence the contribution probability
of the pair is at most \(2Mp_0^2\).

The parameter comparison in
\(\noderef{ParameterHierarchyEventualComparisons}\) gives
\[
M\le \frac{n}{L^2}
\]
in the present range, while \(p_0=\frac12\sqrt{L/n}\).  Consequently
\[
2Mp_0^2\le 2\cdot\frac{n}{L^2}\cdot\frac{L}{4n}
=\frac1{2L}.
\]
Since \(|I|=K\), there are at most
\(C=\noderef{RealChooseTwo}(K)\) unordered pairs in \(I\).  Linearity of
expectation gives
\[
\mathbb E_q Z_{\mathrm{prod}}\le \frac{C}{2L}.
\]

The side conditions also follow from the same range information.  The
inequality \(n_0<n\) gives \(n>0\), and
\(\noderef{FixedSetProductModelConstants}\) gives
\[
c_{\mathrm{prod}}=\frac12 n^{4\varepsilon}>0,
\]
and \(t=C/\sqrt L\ge0\).

For the bounded-difference estimate, change one product coordinate and keep
all others fixed.  Only the external base vertex represented by that coordinate
can change its contribution to \(Z_{\mathrm{prod}}\).  If the vertex is not in
the small class, it contributes nothing.  If it is in the small class, the
definition of \(\noderef{TwoBiteSmallClass}\) bounds the associated fiber size
by
\[
\noderef{TwoBiteSmallCutoff}(\varepsilon,n)=n^{2\varepsilon}.
\]
The cardinality estimate \(\noderef{FinalPairsCardLeHalfSq}\) then bounds the
number of rank-oriented final pairs supported by that one coordinate by
\[
\frac12\left(n^{2\varepsilon}\right)^2=\frac12n^{4\varepsilon}
=c_{\mathrm{prod}}.
\]
The distinct-projection filter can only remove pairs, and dummy or embedded
coordinates change nothing.  Therefore changing one coordinate changes
\(Z_{\mathrm{prod}}\) by at most \(c_{\mathrm{prod}}\).

Finally, since the product distribution of \(Z_{\mathrm{prod}}\) is the
conditional distribution of \(Z_I\) on \(\mathcal C_\pi\) and
\(\mathbb E_qZ_{\mathrm{prod}}\le \mu=C/(2L)\), the conditional event
\[
Z_I\ge \mu+t
\]
is contained in the product-tail event
\[
Z_{\mathrm{prod}}\ge \mathbb E_qZ_{\mathrm{prod}}+t.
\]
Multiplying this conditional inequality by the total mass of
\(\mathcal C_\pi\) gives
\[
\noderef{TwoBiteEventProbability}
  (\mathcal C_\pi\cap\{Z_I\ge\mu+t\})
\le
\Pr_q(Z_{\mathrm{prod}}\ge\mathbb E_qZ_{\mathrm{prod}}+t)\,
\noderef{TwoBiteEventProbability}(\mathcal C_\pi).
\]
The zero-mass case was handled at the start, so all claimed product-model
ingredients hold for the fixed embedding \(\pi\).
\end{proof}
